\documentclass[11pt, a4paper]{article}
\usepackage[a4paper, margin=2cm]{geometry}
\usepackage{fontspec}
\usepackage[english, russian, bidi=basic, provide=*]{babel}
\babelprovide[import, onchar=ids fonts]{russian}
\babelfont{rm}{Noto Sans}
\usepackage{amsmath, amssymb, amsthm}
\usepackage{booktabs}
\usepackage{enumitem}
\usepackage{hyperref}

\title{\textbf{Полные решения заданий}\\ \large Задание 3. Линейные пространства, базис, ранг и собственные значения}
\author{Линейная алгебра}
\date{2024}

\begin{document}

\maketitle

% ==========================================
\section{Раздел 1. Линейная независимость и базис}
% ==========================================

\subsection*{Задание 1.1}
\textbf{Условие:} Проверьте линейную независимость системы векторов:
$$ v_1 = (1, 2, 3), \quad v_2 = (2, 5, 7), \quad v_3 = (3, 7, 10). $$

\textbf{Решение:}
Составим матрицу из векторов-строк и посчитаем ее определитель или приведет к ступенчатому виду:
$$ A = \begin{pmatrix} 1 & 2 & 3 \\ 2 & 5 & 7 \\ 3 & 7 & 10 \end{pmatrix} $$
Вычислим определитель матрицы $A$:
\begin{align*}
\det(A) &= 1 \cdot (5 \cdot 10 - 7 \cdot 7) - 2 \cdot (2 \cdot 10 - 7 \cdot 3) + 3 \cdot (2 \cdot 7 - 5 \cdot 3) \\
&= 1 \cdot (50 - 49) - 2 \cdot (20 - 21) + 3 \cdot (14 - 15) \\
&= 1 \cdot (1) - 2 \cdot (-1) + 3 \cdot (-1) = 1 + 2 - 3 = 0.
\end{align*}
Так как определитель равен нулю, векторы $v_1, v_2, v_3$ являются **линейно зависимыми**. Заметим также очевидное равенство: $v_3 = v_1 + v_2$.

\textbf{Ответ:} Система векторов линейно зависима.

---

\subsection*{Задание 1.2}
\textbf{Условие:} Разложите вектор $x = (5, 9, 14)$ по базису $e_1 = (1, 1, 1)$, $e_2 = (1, 2, 3)$, $e_3 = (2, -1, 1)$.

\textbf{Решение:}
Составим векторное уравнение $c_1 e_1 + c_2 e_2 + c_3 e_3 = x$:
$$ \begin{cases} c_1 + c_2 + 2c_3 = 5 \\ c_1 + 2c_2 - c_3 = 9 \\ c_1 + 3c_2 + c_3 = 14 \end{cases} $$
Решим методом Гаусса. Вычтем 1-е уравнение из 2-го и 3-го:
$$ \begin{cases} c_1 + c_2 + 2c_3 = 5 \\ c_2 - 3c_3 = 4 \\ 2c_2 - c_3 = 9 \end{cases} $$
Умножим второе уравнение на 2 и вычтем из третьего:
$$ (2c_2 - c_3) - 2(c_2 - 3c_3) = 9 - 8 \implies 5c_3 = 1 \implies c_3 = \frac{1}{5} = 0.2. $$
Подставим $c_3$ во 2-е уравнение:
$$ c_2 - 3(0.2) = 4 \implies c_2 = 4.6. $$
Подставим $c_1$ в 1-е уравнение:
$$ c_1 + 4.6 + 2(0.2) = 5 \implies c_1 + 5.0 = 5 \implies c_1 = 0. $$

\textbf{Ответ:} Координаты вектора $x$ в базисе: $x = 0e_1 + 4.6e_2 + 0.2e_3 = (0, \; 4.6, \; 0.2)^T$.

\newpage

% ==========================================
\section{Раздел 2. Ранг матрицы и размерность пространств}
% ==========================================

\subsection*{Задание 2.1}
\textbf{Условие:} Найдите ранг матрицы $A$:
$$ A = \begin{pmatrix} 1 & -2 & 3 & 1 \\ 2 & 1 & -1 & 3 \\ 3 & -1 & 2 & 4 \end{pmatrix} $$

\textbf{Решение:}
Приведём матрицу к ступенчатому виду с помощью элементарных преобразований над строками:
\begin{itemize}
\item $R_2 \leftarrow R_2 - 2R_1$: получаем строку $(0, 5, -7, 1)$
\item $R_3 \leftarrow R_3 - 3R_1$: получаем строку $(0, 5, -7, 1)$
\end{itemize}
$$ \begin{pmatrix} 1 & -2 & 3 & 1 \\ 0 & 5 & -7 & 1 \\ 0 & 5 & -7 & 1 \end{pmatrix} \xrightarrow{R_3 \leftarrow R_3 - R_2} \begin{pmatrix} 1 & -2 & 3 & 1 \\ 0 & 5 & -7 & 1 \\ 0 & 0 & 0 & 0 \end{pmatrix} $$
Количество ненулевых строк в ступенчатом виде равно $2$.

\textbf{Ответ:} $\text{rank}(A) = 2$.

\newpage

% ==========================================
\section{Раздел 3. Собственные значения и собственные векторы}
% ==========================================

\subsection*{Задание 3.1}
\textbf{Условие:} Найдите собственные значения $\lambda$ и собственные векторы матрицы $A = \begin{pmatrix} 4 & 1 \\ 2 & 3 \end{pmatrix}$.

\textbf{Решение:}
1. \textbf{Характеристическое уравнение:} $\det(A - \lambda I) = 0$.
$$ \begin{vmatrix} 4 - \lambda & 1 \\ 2 & 3 - \lambda \end{vmatrix} = 0 $$
$$ (4 - \lambda)(3 - \lambda) - 2 = 0 \implies \lambda^2 - 7\lambda + 12 - 2 = 0 \implies \lambda^2 - 7\lambda + 10 = 0. $$
Корни уравнения по теореме Виета: $\lambda_1 = 5, \quad \lambda_2 = 2$.

2. \textbf{Собственные векторы:}
\begin{itemize}
\item \textbf{Для $\lambda_1 = 5$:}
$$ (A - 5I)v = 0 \implies \begin{pmatrix} -1 & 1 \\ 2 & -2 \end{pmatrix} \begin{pmatrix} v_1 \\ v_2 \end{pmatrix} = \begin{pmatrix} 0 \\ 0 \end{pmatrix} \implies -v_1 + v_2 = 0 \implies v_1 = v_2. $$
Собственный вектор: $v^{(1)} = \begin{pmatrix} 1 \\ 1 \end{pmatrix}$ (или любая его ненулевая пропорция).

\item \textbf{Для $\lambda_2 = 2$:}
$$ (A - 2I)v = 0 \implies \begin{pmatrix} 2 & 1 \\ 2 & 1 \end{pmatrix} \begin{pmatrix} v_1 \\ v_2 \end{pmatrix} = \begin{pmatrix} 0 \\ 0 \end{pmatrix} \implies 2v_1 + v_2 = 0 \implies v_2 = -2v_1. $$
Собственный вектор: $v^{(2)} = \begin{pmatrix} 1 \\ -2 \end{pmatrix}$.
\end{itemize}

\textbf{Ответ:} Собственные значения: $\lambda_1 = 5$, $\lambda_2 = 2$. Собственные векторы: $v^{(1)} = (1, 1)^T$, $v^{(2)} = (1, -2)^T$.

\newpage

% ==========================================
\section{Раздел 4. Бонусные задания и Теория}
% ==========================================

\subsection*{Бонус 1}
\textbf{Условие:} Докажите, что если $\lambda$ --- собственное значение ненулевой обратимой матрицы $A$, то $\frac{1}{\lambda}$ --- собственное значение матрицы $A^{-1}$.

\textbf{Доказательство:}
По определению собственного значения и собственного вектора, существует ненулевой вектор $v \ne 0$ такой, что:
$$ A v = \lambda v. $$
Так как матрица $A$ обратима, ее определитель $\det(A) \ne 0$, а значит, $\lambda \ne 0$.
Умножим обе части равенства слева на $A^{-1}$:
$$ A^{-1} (A v) = A^{-1} (\lambda v) \implies I v = \lambda (A^{-1} v) \implies v = \lambda A^{-1} v. $$
Поделим обе части равенства на скаляр $\lambda \ne 0$:
$$ A^{-1} v = \frac{1}{\lambda} v. $$
Это по определению означает, что $\frac{1}{\lambda}$ является собственным значением матрицы $A^{-1}$ с тем же собственным вектором $v$. \qed

---

\subsection*{Теория}

\textbf{a) Что такое размерность линейного пространства?}
\begin{proof}[Ответ]
Размерность линейного пространства $V$ (обозначается $\dim V$) --- это максимальное число линейно независимых векторов в этом пространстве, либо (эквивалентно) количество векторов в любом его базисе.
\end{proof}

\textbf{b) Какова связь между рангом матрицы и количеством ее линейно независимых строк/столбцов?}
\begin{proof}[Ответ]
Ранг матрицы $\text{rank}(A)$ равен максимальному числу ее линейно независимых строк (или, что то же самое, максимальному числу ее линейно независимых столбцов).
\end{proof}

\end{document}